\section{Arquitetura Proposta}

Esta seção descreve a arquitetura proposta para diagnóstico personalizado de
carteiras. O sistema recebe uma entrada estruturada, executa agentes
especialistas com papéis financeiros distintos, promove um debate
adversarial entre interpretações contrastantes e produz uma saída final
padronizada. A arquitetura não tem por objetivo recomendar compra, venda
ou rebalanceamento: sua função é classificar a situação da carteira em
relação ao perfil do usuário e tornar explícitos os fatores que sustentam
essa classificação. O recorte por diagnóstico, em vez de recomendação, é
deliberado. Ele permite que o sistema admita incerteza, separe risco
observado de lacuna de dados e seja avaliado por métricas que não dependem
de execução em mercado real.

\subsection{Entrada, estado compartilhado e saída}

A entrada do sistema é composta por três grupos de informação. O primeiro
é a carteira, representada por tickers, pesos normalizados e setores. O
segundo é o perfil do usuário, contendo tolerância a risco, horizonte e
objetivo declarado. O terceiro contém dados financeiros externos, quando
disponíveis: preços, fundamentos, notícias e indicadores macroeconômicos.
A separação é intencional. A arquitetura segue operacional mesmo quando
parte dos dados externos não está presente, registrando as lacunas em vez
de inferir o que não foi observado.

Internamente, a execução mantém um estado compartilhado que armazena a
entrada original, as análises produzidas pelos especialistas, as rodadas
de debate, o diagnóstico final e metadados de execução. Cada agente lê o
estado e retorna apenas uma atualização parcial, sem sobrescrever
resultados de outros papéis. A organização também facilita a comparação
com os baselines: o mesmo estado inicial pode ser consumido por uma LLM
única, por um conjunto de especialistas sem debate ou pela arquitetura
completa, mantendo a entrada controlada entre as variantes.

A saída final é um objeto estruturado com cinco campos principais:
classificação discreta, justificativa textual, fatores positivos, fatores
negativos e confiança. As classes possíveis são \textit{estável},
\textit{atenção}, \textit{risco elevado}, \textit{desalinhada} e
\textit{inconclusiva}. A classe \textit{inconclusiva} é central no
problema: ela dá ao sistema um caminho explícito para sinalizar falta de
evidência, em lugar de produzir uma análise aparentemente segura sem base
informacional. Respostas inconclusivas justificadas por ausência de dados
são tratadas, nas métricas, de forma distinta das classificações
afirmativas, para não penalizar o sistema por reconhecer suas próprias
limitações.

\subsection{Agentes especialistas}

Os agentes especialistas decompõem o diagnóstico em perspectivas
financeiras complementares. O agente técnico resume sinais de preço,
médias móveis, volatilidade e tendência, quando esses dados estão
disponíveis. O agente de sentimento interpreta notícias, eventos e
evidências textuais explicitamente fornecidas. O agente fundamentalista
avalia indicadores como valuation, dividendos, margens, lucro e dívida.
O agente macroeconômico resume fatores agregados relevantes para a
carteira, como ciclo de juros e condições setoriais amplas. O agente de
risco avalia concentração, exposição setorial e alinhamento da carteira
ao perfil informado.

A decomposição não assume que cada agente terá evidência suficiente. Uma
regra central dos prompts é distinguir risco observado de lacuna de
dados. Ausência de fundamentos, por exemplo, não deve ser tratada como
fundamento ruim; ela deve ser registrada como incerteza explícita. A
distinção é importante sobretudo na comparação com baselines de agente
único, que tendem a produzir respostas mais assertivas mesmo sob baixa
evidência, escondendo a fragilidade do diagnóstico em texto bem
formulado.

Todos os agentes retornam saídas estruturadas via schemas Pydantic. O
sistema envia o schema esperado no prompt e valida a resposta localmente.
Respostas que não obedecem ao schema são registradas como erro de parsing
e não compõem o diagnóstico final, evitando que saídas malformadas
contaminem as métricas. A taxa de sucesso de parsing é tratada como
métrica operacional explícita e reportada para todos os baselines, como
indicador da confiabilidade do sistema em uso contínuo.

\subsection{Debate adversarial}

O debate adversarial é composto por dois agentes contrastantes. O agente
Bull constrói a leitura favorável mais forte compatível com os dados
disponíveis; o agente Bear, a leitura crítica mais forte. Ambos são
instruídos a não recomendar compra, venda ou rebalanceamento e a não
transformar lacunas de dados em perdas financeiras. O moderador recebe as
duas leituras junto com as análises dos especialistas e produz a síntese
final, atribuindo à classificação um nível de confiança coerente com o
grau de convergência entre as fontes.

O objetivo do debate não é decidir por votação, e sim explicitar
incertezas e fatores contrastantes. Cabe ao Bull organizar evidências
favoráveis e mitigadores; cabe ao Bear tornar visíveis riscos,
fragilidades e pressupostos pouco sustentados. Trata-se, na prática, de
um mecanismo de auditoria interna. Quando Bull e Bear convergem, a
confiança do diagnóstico cresce. Quando divergem fortemente, há razão
para uma classificação intermediária, como \textit{atenção}, ou para uma
classificação \textit{inconclusiva}.

O debate é limitado por um número máximo de rodadas e por uma medida
determinística de convergência lexical entre rodadas, calculada sobre a
interseção dos fatores mencionados por Bull e Bear. A configuração
avaliada utiliza uma rodada por caso, o que mantém a comparação com os
baselines transparente: o custo adicional do debate fica diretamente
refletido no número de chamadas LLM, em tokens e em latência, todos
registrados por execução.

\subsection{Roteamento de modelos}

São avaliadas duas variantes da arquitetura com debate. A variante B4-H
emprega um único modelo em todos os papéis, isolando o efeito da
arquitetura em relação ao modelo subjacente. A variante B4-R utiliza
roteamento heterogêneo, combinando modelos distintos para especialistas,
debate e moderação. O limite de concorrência é configurável e foi ajustado
conforme a estabilidade do provedor. A validação estrutural usou até três
chamadas simultâneas; parte da comparação enriquecida foi reduzida para uma
chamada por vez após respostas temporárias de sobrecarga do serviço. Essa
diferença é registrada nos logs e entra nas métricas operacionais de
latência e chamadas, sem alterar a entrada compartilhada entre baselines.

Na configuração avaliada, B4-H utiliza \texttt{deepseek-v4-pro} em todos
os papéis. Em B4-R, os agentes técnico e de risco utilizam
\texttt{nemotron-3-super}; sentimento, fundamentos, macroeconomia e
moderação utilizam \texttt{deepseek-v4-pro}; Bull e Bear utilizam
\texttt{minimax-m2.7}. As páginas oficiais do Ollama classificam
\texttt{deepseek-v4-pro} como uso ``extra high'', e
\texttt{minimax-m2.7} e \texttt{nemotron-3-super} como uso ``medium''
\citep{ollamadeepseek2026,ollamaminimax2026,ollamanemotron2026}. A escolha
de modelos por papel não é apresentada como ótima em sentido absoluto. Ela
operacionaliza a hipótese de que papéis diferentes podem se beneficiar de
modelos com perfis distintos, enquanto B4-H atua como controle homogêneo
dessa hipótese.

\subsection{Implementação e rastreabilidade}

A arquitetura é implementada em Python, com LangGraph para orquestração
do estado compartilhado entre agentes e Pydantic para os schemas de
entrada e saída. As chamadas aos modelos passam por uma camada uniforme,
o que permite trocar provedores sem alterar a lógica dos agentes. Cada
chamada LLM passa por um envelope que registra o agente de origem, o
baseline em execução, o caso analisado, o modelo utilizado e o tempo de
resposta.

A execução completa registra, em SQLite, informações detalhadas por
caso, baseline, agente e chamada LLM: modelo, prompt enviado, resposta
bruta, resposta parseada, status, latência, tokens quando disponíveis e
erro quando aplicável. Os registros sustentam tanto a depuração quanto
as métricas reportadas, e permitem reexecutar análises agregadas sem
acionar novamente os modelos. Em outras palavras, é possível auditar de
forma objetiva se diferenças entre baselines vêm de mudanças na
arquitetura, do modelo utilizado ou de comportamento intermitente do
provedor.
